% ============================================================================
%  SECCIÓN 1.6: DEFINICIÓN DE VARIABLES
% ============================================================================

\section{Definici\'on de variables}

La concreci\'on de los aspectos de la realidad o del problema a los que se
dirige la atenci\'on se realiza definiendo las variables. Una variable es un
aspecto de la realidad que cambia en distintas situaciones y para diferentes
objetivos y sujetos; las variables permiten agrupar, diferenciar, ordenar,
distribuir y relacionar los elementos de la realidad \citep{hernandez2014}.
Las variables no est\'an dadas de antemano, sino que hay que elaborarlas,
siendo en un principio conceptos generales que se concretan para darles
significado.

La variable es una determinada caracter\'istica o propiedad del objeto de
estudio a la cual se observa y/o cuantifica en la investigaci\'on. Para el
presente proyecto se han identificado las variables que se presentan en la
siguiente tabla, con su definici\'on conceptual, sus dimensiones, sus
indicadores y los instrumentos de recolecci\'on de datos que permitir\'an
medirlas.

\begin{table}[H]
  \centering
  \caption{Operacionalizaci\'on de las variables de la investigaci\'on}
  \label{tab:variables}
  \contenidoConFuente{%
    \begin{tablaAPA}
      \begin{tabular}{@{}p{2.3cm}p{3.5cm}p{2.8cm}p{3.8cm}p{2.3cm}@{}}
        \toprule
        \textbf{Variable} & \textbf{Definici\'on conceptual} &
        \textbf{Dimensiones} & \textbf{Indicadores} &
        \textbf{Instrumentos} \\
        \midrule
        Acceso a la informaci\'on deportiva &
        Facilidad con la que el p\'ublico obtiene informaci\'on sobre la oferta
        deportiva disponible &
        Nivel de dificultad; tiempo de b\'usqueda; canales utilizados &
        I1. Porcentaje de personas que desconoce la oferta deportiva de su
        ciudad; I2. tiempo promedio de b\'usqueda &
        Encuesta \\
        \addlinespace
        Visualizaci\'on interactiva &
        Representaci\'on gr\'afica e interactiva de la informaci\'on de los
        escenarios &
        Usabilidad del mapa; utilidad de los filtros &
        I1. Usabilidad percibida (escala); I2. tasa de uso del mapa
        interactivo &
        Encuesta, gu\'ia de observaci\'on \\
        \addlinespace
        Cobertura del sistema &
        N\'umero y distribuci\'on de los escenarios registrados en la plataforma &
        Cantidad de recintos; distribuci\'on geogr\'afica &
        I1. Escenarios registrados (total); I2. departamentos cubiertos &
        Revisi\'on documental, base de datos del sistema \\
        \bottomrule
      \end{tabular}
    \end{tablaAPA}%
  }{Elaboraci\'on propia en base a \citet{hernandez2014}.}
\end{table}
